% Software architecture, drawn at the manuscript text width.
% Shared figure styles are supplied by paper_layout.tex.
\begingroup
\colorlet{scInk}{black}
\colorlet{scInkM}{black!65}
\colorlet{scSpF}{white}
\colorlet{scRxF}{white}
\colorlet{scRefF}{white}
\colorlet{scHiF}{white}
\colorlet{scAccentF}{white}
\colorlet{scAlertF}{white}
\colorlet{scMarkF}{white}
\colorlet{scSpE}{black}
\colorlet{scRxE}{black}
\colorlet{scReac}{black}
\colorlet{scProd}{black}
\colorlet{scAccent}{black}
\colorlet{scRef}{black}
\colorlet{scMark}{black!65}
% Panel letters: the shared plain "(a)" of paper_layout.tex.

\tikzset{
  minisp/.style={circle,fill=white,draw=scSpE,line width=0.45pt,
    minimum size=3.2mm,inner sep=0pt},
  minirx/.style={rectangle,fill=white,draw=scRxE,line width=0.45pt,
    minimum size=2.7mm,inner sep=0pt},
  % the route arrows between panels
  pipe/.style={-{Stealth[length=1.6mm,width=1.2mm]},draw=black!55,
    line width=0.55pt},
  % one row of an adapter list: icon column, then name, then format
  rowname/.style={font=\scBody,text=scInk,anchor=base west,
    execute at begin node={\bfseries}},
  rowgloss/.style={font=\scSmall,text=scInkM,anchor=base west},
  % a row whose name is optically centred on its icon (no second line)
  rowmid/.style={font=\scBody,text=scInk,anchor=west,
    execute at begin node={\bfseries}},
}

% Column geometry, in cm, on the 16.40 canvas.
\def\W{16.40}
\def\ax{0}\def\aw{4.05}
\def\bx{4.55}\def\bw{7.30}
\def\cx{12.35}\def\cw{4.05}
\def\rowh{3.40}          % height of the top row of panels
\def\dy{-3.82}\def\dh{2.55}  % analysis band
% Icon column: the widest glyph reaches 0.62 cm from its centre, so names start
% 0.78 cm out.  At the previous 0.62/1.06 the hypergraph ran 1.8 mm into its name.
\def\icon{0.80}\def\name{1.58}
% Row centres, shared by panels (a) and (c) so the two columns align.
\def\rya{-1.18}\def\ryb{-2.05}\def\ryc{-2.92}

\begin{tikzpicture}[syn]
\path[use as bounding box] (0,0) rectangle (\W,-6.37);

% =========================================================================
% (a) input adapters
% =========================================================================
\scHead{(\ax+0.22,-0.16)}{a}{Input adapters}{}

% --- Query: a database cylinder
\begin{scope}[shift={(\ax+\icon,\rya+0.16)},draw=scAccent,line width=0.5pt]
  \draw[fill=scAccentF] (0,0) ellipse (0.17 and 0.065);
  \draw[fill=scAccentF] (-0.17,0) -- (-0.17,-0.26)
    arc[start angle=180,end angle=360,x radius=0.17,y radius=0.065] -- (0.17,0);
  \draw (-0.17,-0.13) arc[start angle=180,end angle=360,x radius=0.17,y radius=0.065];
\end{scope}
\node[rowname]  at (\ax+\name,\rya+0.04) {Query};
\node[rowgloss] at (\ax+\name,\rya-0.22) {KEGG modules};

% --- IO: an XML tag
\node[fill=white,draw=scRef,line width=0.4pt,
      font=\scSmall,text=scRef,inner xsep=2.5pt,inner ysep=1.5pt]
  at (\ax+\icon,\ryb) {<xml>};
\node[rowname]  at (\ax+\name,\ryb+0.04) {IO};
\node[rowgloss] at (\ax+\name,\ryb-0.22) {SBML L3 Core};

% --- Construct: a rule firing
\begin{scope}[shift={(\ax+\icon,\ryc)}]
  \node[minisp] (ia) at (-0.42,0) {};
  \node[minirx,fill=scMarkF,draw=scMark] (ir) at (0,0) {};
  \node[minisp] (ib) at (0.42,0) {};
  \draw[reac] (ia)--(ir); \draw[prod] (ir)--(ib);
\end{scope}
\node[rowname]  at (\ax+\name,\ryc+0.04) {Construct};
\node[rowgloss] at (\ax+\name,\ryc-0.22) {SMARTS / DPO rules};

% =========================================================================
% (b) the canonical object
% =========================================================================
\draw[pipe] (\ax+\aw+0.10,-1.70) -- (\bx-0.10,-1.70);

\scHead{(\bx+0.22,-0.16)}{b}{Reaction records}{}
\node[anchor=north east,font=\scBody,text=scInk] at (\bx+\bw-0.22,-0.155)
  {\textbf{SynCRN}};

% One reaction with two substrates, one with a catalyst that is consumed and
% returned.  Both are exactly what a net-stoichiometry view would lose.
\begin{scope}[shift={(\bx+0.95,-1.85)}]
  \node[minisp] (a)  at (0,0.58) {};
  \node[minisp] (b)  at (0,-0.58) {};
  \node[minirx] (r1) at (1.35,0) {};
  \node[minisp] (c)  at (2.70,0) {};
  \node[minisp] (d)  at (2.70,-0.92) {};
  \node[minirx] (r2) at (4.05,-0.38) {};
  \node[minisp] (e)  at (5.40,0.14) {};
  \node[minisp] (f)  at (5.40,-0.88) {};
  \node[minisp,fill=scHiF,draw=scHiE] (cat) at (4.05,0.78) {};
  \draw[reac] (a)--(r1); \draw[reac] (b)--(r1);
  \draw[prod] (r1)--(c);
  \draw[reac] (c)--(r2); \draw[reac] (d)--(r2);
  \draw[prod] (r2)--(e); \draw[prod] (r2)--(f);
  % the catalyst is consumed and returned by the same reaction
  \draw[cat] (cat) to[out=200,in=110] (r2);
  \draw[cat] (r2) to[out=70,in=-20] (cat);
  \node[stoi,fill=white,inner sep=1pt] at (0.46,0.33) {2};
  \node[font=\scSmall,text=scHiE,anchor=west] at (4.35,0.78) {catalyst};
\end{scope}

\node[anchor=north,font=\scBody,text=scInkM] at (\bx+0.5*\bw,-3.02)
  {$\mathcal N=(X,\mathscr R,\mathbf S^{-},\mathbf S^{+})$};

% =========================================================================
% (c) derived views
% =========================================================================
\draw[pipe] (\bx+\bw+0.10,-1.70) -- (\cx-0.10,-1.70);

\scHead{(\cx+0.22,-0.16)}{c}{Derived representations}{}

% --- stoichiometry: an illustrative matrix
\begin{scope}[shift={(\cx+\icon,\rya)}]
  \node[font=\scBody,text=black,inner sep=0pt] at (0,0)
    {$\left(\begin{smallmatrix}-1&0\\1&-1\\0&1\end{smallmatrix}\right)$};
\end{scope}
\node[rowmid] at (\cx+\name,\rya) {stoichiometry};

% --- Petri net: a marked place firing
\begin{scope}[shift={(\cx+\icon,\ryb)}]
  \node[place,minimum size=3.0mm] (p1) at (-0.42,0) {};
  \node[token] at (-0.42,0) {};
  \node[trans,minimum height=3.0mm] (pt) at (0,0) {};
  \node[place,minimum size=3.0mm] (p2) at (0.42,0) {};
  \draw[reac] (p1)--(pt); \draw[prod] (pt)--(p2);
\end{scope}
\node[rowmid] at (\cx+\name,\ryb) {Petri net};

% --- hypergraph: a two-tail edge
\begin{scope}[shift={(\cx+\icon,\ryc)}]
  \node[minisp] (h1) at (-0.46,0.20) {};
  \node[minisp] (h2) at (-0.46,-0.20) {};
\node[draw=scRxE,fill=white,minimum width=1.8mm,
        minimum height=5.0mm,inner sep=0pt] (hr) at (0,0) {};
  \node[minisp] (h3) at (0.46,0) {};
  \draw[reac] (h1)--(hr); \draw[reac] (h2)--(hr); \draw[prod] (hr)--(h3);
\end{scope}
\node[rowmid] at (\cx+\name,\ryc) {hypergraph};

% =========================================================================
% (d) structural analyses
% =========================================================================
\scHead{(0.22,\dy-0.16)}{d}{Structural analyses}{}

\def\dcw{3.825}\def\dcy{\dy-0.60}\def\dch{1.80}
% card centres, precomputed: TikZ would read the parentheses of an inline
% expression as a nested coordinate.
\pgfmathsetmacro{\cdA}{0.22+0.5*\dcw}
\pgfmathsetmacro{\cdB}{0.22+1*(\dcw+0.22)+0.5*\dcw}
\pgfmathsetmacro{\cdC}{0.22+2*(\dcw+0.22)+0.5*\dcw}
\pgfmathsetmacro{\cdD}{0.22+3*(\dcw+0.22)+0.5*\dcw}
\foreach \i/\ttl in {0/{CRNT structure},1/{conservation},
                     2/{Petri structure},3/{realizability}} {
  \pgfmathsetmacro{\cxx}{0.22+\i*(\dcw+0.22)}
  \node[anchor=north west,font=\scBody,text=scInk]
    at (\cxx+0.16,\dcy-0.115) {\bfseries\ttl};
}

% --- card 1: the deficiency identity
\node[font=\fontsize{9.5}{11}\selectfont,text=scInk] at (\cdA,\dy-1.42)
  {$\delta=n-\ell-\operatorname{rank}\mathbf S$};
\node[font=\scSmall,text=scInkM,align=center,text width=3.45cm]
  at (\cdA,\dy-2.02) {exact rational rank};

% --- card 2: a conservation law spans the chain
\begin{scope}[shift={(\cdB,\dy-1.36)}]
  \node[minisp] (c1) at (-1.30,0) {};
  \node[minirx] (cr) at (-0.65,0) {};
  \node[minisp] (c2) at (0,0) {};
  \node[minirx] (cs) at (0.65,0) {};
  \node[minisp] (c3) at (1.30,0) {};
  \draw[reac] (c1)--(cr); \draw[prod] (cr)--(c2);
  \draw[reac] (c2)--(cs); \draw[prod] (cs)--(c3);
  \draw[decorate,decoration={brace,amplitude=2.4pt,raise=1.6pt},
        draw=scAccent,line width=0.4pt] (1.44,-0.19) -- (-1.44,-0.19);
\end{scope}
\node[font=\scSmall,text=scInkM,align=center,text width=3.45cm]
  at (\cdB,\dy-2.02) {minimal $y\geq0$ with $y^{\mathsf T}\mathbf S=0$};

% --- card 3: a siphon --- once empty, never refilled
\begin{scope}[shift={(\cdC,\dy-1.36)}]
  \node[minisp,fill=scAlertF,draw=scAlert] (s1) at (-1.20,0) {};
  \node[minirx] (sr) at (-0.60,0) {};
  \node[minisp] (s2) at (0,0) {};
  \node[minirx] (ss) at (0.60,0) {};
  \node[minisp] (s3) at (1.20,0) {};
  \draw[reac] (s1)--(sr); \draw[prod] (sr)--(s2);
  \draw[reac] (s2)--(ss); \draw[prod] (ss)--(s3);
  \draw[scAlert,line width=0.55pt] (-1.32,-0.12)--(-1.08,0.12);
\end{scope}
\node[font=\scSmall,text=scInkM,align=center,text width=3.45cm]
  at (\cdC,\dy-2.02) {minimal siphons and traps};

% --- card 4: a replayable firing sequence
\begin{scope}[shift={(\cdD,\dy-1.30)}]
  \foreach \x/\n in {-1.15/1,-0.575/1,0/2,0.575/1,1.15/2} {
    \node[draw=scSpE,fill=white,minimum width=3.8mm,
          minimum height=2.3mm,inner sep=0pt] at (\x,0) {};
    \ifnum\n=1 \node[token] at (\x,0) {};
    \else \node[token] at (\x-0.06,0) {}; \node[token] at (\x+0.06,0) {}; \fi
  }
  \draw[-{Stealth[round,length=1.6mm,width=1.3mm]},draw=scMark,line width=0.5pt]
    (-1.38,-0.30)--(1.38,-0.30);
\end{scope}
\node[font=\scSmall,text=scInkM,align=center,text width=3.45cm]
  at (\cdD,\dy-2.02) {a firing certificate, or inconclusive};

\end{tikzpicture}
\endgroup
